\chapter{绪\quad 论}

\section{课题背景及研究的目的和意义}

能源系统的高效运行依赖对设备状态、用户负荷、环境扰动和控制策略的综合判断。随着传感器网络和运行管理平台逐步普及，研究者能够获得比传统人工巡检更细致的时序数据，但这些数据也带来了缺失、异步、异常值和解释困难等新问题。博士论文模板需要容纳较长的理论论证、模型推导、实验表格和工程讨论，因此示例正文采用能源系统这一公开主题，模拟正式论文中常见的章节层次。

本文示例的目的有两项。其一，给出一份可以直接编译的哈尔滨工业大学博士论文 LaTeX 脚手架，覆盖官方书写范例中的主要页面。其二，用一组公开、无隐私的研究叙述验证封面、目录、章节、图表、公式和后置材料是否能在统一构建链路中稳定生成。

\section{数据驱动能源系统研究概况}

数据驱动方法通常需要在工程约束和统计模型之间建立稳定接口。对于能源站运行问题，原始数据可以分为三类：设备侧数据、环境侧数据和调度侧数据。设备侧数据描述温度、压力、流量、功率和启停状态；环境侧数据描述室外温度、湿度和辐照度；调度侧数据则记录负荷需求、价格信号和人工策略。

\subsection{能源系统状态建模}

状态建模的关键在于选择合适的时间尺度。过短的窗口容易放大噪声，过长的窗口又可能掩盖设备切换和负荷突变。示例模板通过正文段落、编号小节和公式排版展示博士论文常见写法：
\begin{equation}
  \eta_t = \frac{Q_t}{E_t + \epsilon},
\end{equation}
其中，$\eta_t$ 表示时刻 $t$ 的综合能效，$Q_t$ 表示有效供能量，$E_t$ 表示输入能耗，$\epsilon$ 为避免除零的稳定项。

\subsection{调度优化问题}

调度优化需要同时处理能耗、舒适性和设备寿命。一个简化目标函数可写为
\begin{equation}
  \min \sum_{t=1}^{T}\left(\alpha E_t + \beta P_t^{\max} + \gamma R_t\right),
\end{equation}
其中 $P_t^{\max}$ 为峰值负荷惩罚，$R_t$ 为运行风险项，$\alpha$、$\beta$、$\gamma$ 为权重。

\section{本文的主要研究内容}

本文示例围绕多源数据建模与调度方法展开。第 1 章说明研究背景和论文结构；第 2 章给出多源数据建模框架；第 3 章展示滚动调度流程和模板验证结果。后置部分依次给出结论、参考文献、创新成果、评阅与答辩信息、原创性声明、致谢和个人简历。
